\newpage
\section{RAG na regulaminie forum (5.0)}
\label{sec:rag}

Moderator, który usuwa wpis, powinien wskazać złamany punkt regulaminu. Model sam z siebie nie zna regulaminu tego
konkretnego forum, więc w trybie RAG dostaje jego fragmenty w prompcie.

\subsection{Baza wiedzy}

Regulamin \kod{data/rules/forum\_rules.md} ma 6 rozdziałów i 14 paragrafów. Plik \kod{rules.py} dzieli go na paragrafy
po nagłówkach \kod{\#\#\# §X.Y}; każdy paragraf to jeden dokument. Rozdziały §1 (postanowienia ogólne) i §6 (zasady
moderacji) nie opisują zachowań użytkowników, więc nie biorą udziału w wyszukiwaniu (\kod{skip\_sections} w
\kod{config.yaml}). Regulamin jest po polsku, ale część komentarzy jest po angielsku, dlatego na końcu paragrafów
dopisałem w nawiasie krótki angielski opis.

\subsection{Wyszukiwanie}

\kod{retrieval.py} ma trzy metody. BM25 działa na słowach po prostym ujednoliceniu: małe litery, usunięte polskie znaki
i obcięcie słowa do 5 pierwszych liter (prymitywny stemming, który łączy np. ,,groźby'' i ,,grożenie''). Wyszukiwanie
wektorowe używa \kod{nomic-embed-text} z prefiksami \kod{search\_document:} i \kod{search\_query:}, których ten model
wymaga. Metoda hybrydowa łączy oba rankingi przez RRF.

\zrzut{code_retrieval}{Kod: wyszukiwarka paragrafów (\kod{retrieval.py}).}

\subsection{Ocena z cytatem}

Do promptu trafiają 3 najlepsze paragrafy oraz zawsze §2.3 (dozwolona krytyka), żeby przy zwykłym komentarzu model
miał paragraf, na który może się powołać. W schemacie JSON pole \kod{rule} ma listę dozwolonych wartości
(\kod{enum}) równą numerom znalezionych paragrafów, więc model nie może wymyślić paragrafu. Na koniec kod sprawdza
odpowiedź jeszcze raz: paragraf spoza regulaminu albo usunięcie bez paragrafu zmienia decyzję na \kod{review}.

\zrzut{code_rag}{Kod: ocena komentarza z regulaminem (\kod{rag.py}).}

Pierwsza wersja nie miała listy \kod{enum}, a w instrukcji było ,,podaj paragraf albo null''. Model bardzo często
wybierał \kod{null}, także przy oczywistej groźbie, i takie komentarze trafiały do \kod{review} zamiast do
\kod{remove}. Ograniczenie pola w schemacie rozwiązało ten problem -- wyniki obu wersji są w rozdziale
\ref{sec:ewaluacja}.

\subsection{Przykłady}

\zrzut{rag_spam}{Wynik: \kod{moderator rag} dla reklamy z kodem rabatowym.}

\zrzut{rag_threat}{Wynik: \kod{moderator rag} dla groźby.}
